% Transcription — fonds Alexandre Grothendieck, shelfmark 80, batch 3 (pages 41-60).
% Established from the facsimile archives/batches/80/batch-03.pdf, cut from a
% copy of 80.pdf fetched through the project's own relay
% (https://grothendieck-relay.onrender.com/source/80.pdf), not uploaded by the
% user; PDF page k = archivists' page 40+k, no cover sheet. Per the skill
% transcribe-grothendieck. The folder is seventy-two pages in four batches
% (1-20, 21-40, 41-60, 61-72), transcribed in parallel; this is the third.
% Batch 1 was read for notation before this pass; batch 2
% (transcripts/80/batch-02.fr.tex) appeared while this pass was being written
% and was read at the end, and its notation followed (\mathcal{X} for his
% cursive barred X, \mathfrak{P}_3 for his gothic P); neither was edited.
%
% Pass: Opus 5.5 (claude-opus-5-5), 2026-09-26 — first pass, unchecked against the pages by a human.
%
% Dating and title. \dating is the folder's own date, copied verbatim from
% src/content/catalogue.ts; the group « Géométrie et topologie combinatoire »
% (69-89) is not added, since the folder has a date of its own. Title from the
% catalogue, trailing full stop dropped (as batch 1, per the coordinator).
%
% Pages skipped: 55 and 59, blank cream separator sheets with nothing on them.
% Pages summarised: none. Pages given as a figure described in one French
% \note, with their labels and counts set: 56, 57, 58 (small landscape cards,
% ink drawings of line arrangements). Pages transcribed from his hand: 41-54
% (brown ink, one continuous run) and 60 (pencil, a fresh start).
%
% His pagination: pages 41-54 continue batch 2's run (numbered 1-5 on the
% rectos of pp. 31-39), numbered on the rectos only, top left: 6 (p. 41),
% 7 (p. 43), 8 (p. 45), 9 (p. 47), 10 (p. 49), 11 (p. 51), 12 (p. 53); the
% even pages are the versos and carry the text on. Recorded in a \note at each
% numbered page. His equation numbers continue batch 2's (1)-(18): here
% (19)-(34), and his conditions (A1)-(A6). Page 60 is unnumbered and begins a
% new text, « Système de Pseudodroites », which runs into batch 4.
%
% What the batch is about.
%   41-49  Proof strategy for batch 2's « Th. » (the functor from maps
%          X that are pseudo-line arrangements to oriented 1-dimensional
%          arrangements (K, (D_i), omega) is fully faithful): reconstruct X
%          from (K, (D_i), omega) by building, for each vertex s, the pair
%          Z_s (19) of opposite circular orders on the oriented edges at s,
%          and for each oriented edge a the transport psi_a : Z_s -> Z_s'
%          (20)-(21). Conditions: (22) rho^{nu_s}(D_i) = -D_i; the
%          Euler-Poincaré condition (23) c_0 - c_1 + c_2 = 1 for X to be the
%          projective plane. Construction of Z_s: (24)-(25) the polygonal
%          structure on I_s induced by a line Delta not through s via
%          phi_{s,Delta} : I_s -> S_Delta; (A1) independence of Delta; (26)
%          the double cover vec I_s -> I_s; (27); counting 2^{nu_s - 1}
%          choices (28); the triangle rule (29) fixing the lift u_{s,Delta};
%          (A2) circularity, (30), (A3) independence of Delta'; (31) the map
%          vec I \ vec I_s -> vec Z_s; transport along an edge (32), (A4);
%          the special case where no line misses both ends of the edge (33),
%          psi_a = psi_b' psi_b (34), and the doubt about the case of n
%          concurrent lines cut by one transversal.
%   50-54  The « canonical » orientation of a triangle K_J, J in P_3(I):
%          (1er cas) a fourth line missing its vertices, and a Proposition
%          (simple arrangements of 4 lines -> 1-dimensional ones is an
%          equivalence), (A5); (2ème cas) every line through a vertex: a
%          six-line configuration a_i, alpha_i, a'_i, cases I and II, (A6);
%          the residual standard arrangement of 6 lines with 4 triple
%          points; « En résumé, Théorème »: the functor from arrangements
%          without bigons to 1-dimensional ones without bigons is fully
%          faithful; (A1)-(A6) should determine the essential image.
%   56-58  Drawings: arrangements with regions labelled by their number of
%          sides and « id »; region counts 31 = 25 + 5 + 1 (p. 56), and
%          31 = 5 + 15 + 10 + 1 with types 1_6 icosaèdre, 5_1 hexagone,
%          5_4, 5_2 triangle bordé, 5_5, 10_3 carré (p. 58).
%   60     « Système de Pseudodroites »: definition as a map on the real
%          projective plane with even vertex orders >= 4, pseudo-lines as
%          classes of edges meeting pairwise once; first consequences
%          (non-contractible, realisable by straight lines for n <= 6 in
%          general position); the orientation double cover X~ ≃ S^2, the two
%          hemispheres of X~ \ D~ and the orientations of D.
%
% Notation fixed once for the batch: \mathcal{X} for his cursive barred X
% (the map), as batch 2; \mathfrak{Z} for his three-shaped Z (the sets of
% circular orders at a vertex), with \vec as he draws it; \mathfrak{Z}^{1}
% for its variant with a small raised mark (reading of the mark doubtful,
% said in a note); \mathcal{C}_E for his curly C (the set of circular orders
% on E); \nu_s for the number of lines through s; \mathfrak{P}_3(I) for his
% gothic P; \psi for his psi, also when it shrinks to a « v » on p. 49.
%
% Readings to check first: the superscript of \mathfrak{Z}^{1}; the sign
% between \vec{\mathfrak{Z}}_s and \mathcal{C} in (19); the diagonal marginal
% notes on pp. 42, 43, 46, 48, 49, 51, 53 (mostly \ill{}); the last lines of
% p. 48; the six-line argument of pp. 51-53; the word « concourantes » on
% p. 50, where a triangle asks for non-concurrent lines.

\documentclass[11pt,a4paper]{article}
% Le préambule commun à toutes les transcriptions du fonds Grothendieck.
%
% Il tient en une idée : **l'incertitude se compose, elle ne se résout pas.**
% Une transcription qui lisse un mot douteux a détruit la seule chose pour
% laquelle on la faisait — un lecteur doit pouvoir savoir, à chaque endroit,
% ce qui était sur le feuillet et ce qui a été ajouté. Les six macros qui
% suivent sont donc l'appareil critique, et non de la décoration.
%
% Les mêmes commandes sont reconnues par `scripts/render.mjs`, qui produit la
% vue de lecture du site. Une macro ajoutée ici doit l'être aussi là-bas,
% faute de quoi le rendu échouera bruyamment — ce qui est voulu : un rendu
% silencieusement faux est pire qu'un rendu absent.

\ProvidesPackage{grothendieck}[2026/08/08 Transcriptions du fonds Grothendieck]

\usepackage{amsmath,amssymb,amsthm}
\usepackage{mathrsfs}
\usepackage{stmaryrd}
% Les diagrammes commutatifs sont partout dans ce fonds : c'est la langue
% même de la géométrie algébrique de Grothendieck, et une transcription qui
% les rendrait en prose ne transcrirait rien.
\usepackage{tikz-cd}
% Français par défaut — c'est la langue des feuillets — mais l'anglais est
% chargé aussi : la traduction et le résumé sont des documents anglais, et la
% typographie française (espaces avant les deux-points) y serait une faute.
% Ces documents basculent par \selectlanguage{english} en tête de corps.
\usepackage[english,french]{babel}
\usepackage{fontspec}
\usepackage{geometry}
\geometry{a4paper,margin=25mm,marginparwidth=28mm}
\usepackage{marginnote}
\usepackage{xcolor}
% ulem, not soul. `soul`'s \st and \ul are built on a character-by-character
% scan that XeLaTeX's font machinery breaks: the document fails with an
% undefined control sequence rather than degrading. ulem does the same two jobs
% and is Unicode-engine safe. `normalem` stops it hijacking \emph, which the
% transcriptions use for Grothendieck's own emphasis.
\usepackage[normalem]{ulem}
\usepackage{hyperref}
\hypersetup{colorlinks=true,linkcolor=black,urlcolor=black,pdfborder={0 0 0}}

% --- Métadonnées du lot -----------------------------------------------------
% Reprises telles quelles par le script de rendu, qui les affiche en tête de
% la vue de lecture. Elles fixent ce que le fichier prétend couvrir : sans
% elles, une transcription flotte sans point d'attache dans un fonds de seize
% mille feuillets.

\newcommand{\folder}[1]{\gdef\grothFolder{#1}}
\newcommand{\batch}[1]{\gdef\grothBatch{#1}}
\newcommand{\leaves}[2]{\gdef\grothFirst{#1}\gdef\grothLast{#2}}
\newcommand{\foldertitle}[1]{\gdef\grothFoldertitle{#1}}
\gdef\grothFolder{?}\gdef\grothBatch{?}\gdef\grothFirst{?}\gdef\grothLast{?}\gdef\grothFoldertitle{}

% --- Le résumé --------------------------------------------------------------
%
% Il ouvre la lecture modernisée, sous le titre « Résumé », et il est composé
% en petit. Ce n'est pas un ornement : c'est le seul passage du document qui
% ne soit pas des mathématiques, et un lecteur doit pouvoir le reconnaître
% avant de l'avoir lu — puis le sauter s'il connaît déjà le sujet, ou s'y
% arrêter s'il ne le connaît pas. Le filet qui le suit marque la reprise du
% corps.

\newenvironment{resume}
  {\small\setlength{\parskip}{0.45em}}
  {\par\medskip\hrule\medskip}

% --- Mots-clés ---------------------------------------------------------------
%
% En anglais, en clôture du résumé de la lecture modernisée : le vocabulaire
% moderne sous lequel chercher ce que les feuillets construisent. Le manifeste
% du site (`npm run manifest`) les extrait pour étiqueter la cote entière —
% cette ligne est la source unique des tags, il n'y a pas de fichier à côté.
\newcommand{\keywords}[1]{\par\smallskip\noindent
  {\footnotesize\textsc{Keywords} --- \textit{#1}}\par}

% --- L'appareil critique ----------------------------------------------------

% Le numéro de feuillet, en marge. C'est l'ancre qui permet de régler un
% désaccord sur une formule en regardant *un* feuillet plutôt qu'une chemise
% de six cents. Le site s'en sert aussi pour tourner les pages du fac-similé
% quand on fait défiler la transcription.
\newcommand{\leaf}[1]{%
  \marginnote{\footnotesize\color{gray!70}\textsf{#1}}%
  \label{leaf:#1}\ignorespaces}

% Illisible : on ne devine pas. Un mot inventé qui se lit comme les autres est
% le pire résultat possible.
\newcommand{\ill}{\textcolor{red!60!black}{[\,\ldots\,]}}

% Lecture douteuse : le mot est donné, et signalé comme incertain.
\newcommand{\uncertain}[1]{\uline{#1}}

% Ajout de l'éditeur — une lettre manquante, un mot sous-entendu.
\newcommand{\add}[1]{\textcolor{blue!50!black}{[#1]}}

% Biffé par Grothendieck lui-même. Ce qu'il a rayé fait partie du document :
% une rature dit souvent où la pensée a changé de direction.
\newcommand{\struck}[1]{\sout{#1}}

% Note du transcripteur, sur l'état matériel du feuillet.
\newcommand{\note}[1]{\par\smallskip{\small\color{gray!60!black}\textit{[#1]}}\par\smallskip}

% Note marginale de Grothendieck, distincte du corps du texte.
\newcommand{\marginal}[1]{\par\smallskip\noindent\hspace{1em}%
  {\small\color{gray!60!black}#1}\par\smallskip}

% --- Titre ------------------------------------------------------------------

\AtBeginDocument{%
  \begin{center}
    {\large\bfseries Fonds Alexandre Grothendieck --- cote n\textsuperscript{o}~\grothFolder}\\[2pt]
    {\itshape\grothFoldertitle}\\[4pt]
    {\small feuillets \grothFirst--\grothLast\ (lot \grothBatch)}\\[2pt]
    {\footnotesize\color{gray!60!black}Transcription établie d'après le fac-similé numérisé
     par l'Université de Montpellier.\\
     Les lectures incertaines sont soulignées, les illisibles notées [\,\ldots\,],
     les ajouts entre crochets.}
  \end{center}
  \vspace{1em}}

\endinput


\folder{80}
\batch{3}
\pages{41}{60}
\dating{[à partir de 1981]}
\watermark{Édition de démonstration}
\foldertitle{Systèmes de pseudo-droites : notes manuscrites (s.d.)}

\begin{document}

\page{41}
\note{sa page \emph{6}, en haut à gauche ; la suite numérotée commencée page 31 (lot 2) continue sur les rectos, les pages paires étant les versos. Le texte poursuit le « Th. » de la page 39 : le foncteur $\mathcal{X} \mapsto (K, (D_i), \omega)$ est pleinement fidèle}
Pour le voir, on va montrer comment la carte $\mathcal{X}$ peut être reconstruite à partir de l'arrangement $1$-dimensionnel orienté. Par le théorème général des cartes, il \uncertain{revient au même}, quand les sommets de $K$ sont d'ordre $\geqslant 4$ (a fortiori $\geqslant 3$), de construire, à partir de $(K, (D_i)_{i \in I}, \omega)$, :

1°) Pour \struck{tout} $s \in S$ (ens.\ des sommets de $K$), l'ens
\[
(19) \qquad \vec{\mathfrak{Z}}_s \subset \mathcal{C}_{\vec{A}_s}
\]
des deux ordres circulaires (inverses l'un de l'autre) canoniques sur l'ens $\vec{A}_s$ des arêtes \struck{\ill{}} de $K$ d'origine $s$, qu'on peut identifier à : l'ens $\vec{I}_s$ des ps.\ droites orientées $\vec{D}_i$ passant par $s$.
\note{la formule est encadrée. Le signe entre $\vec{\mathfrak{Z}}_s$ et $\mathcal{C}$ est \uncertain{lu} $\subset$ (on pourrait lire $\in$) ; il est suivi d'un « $A$ » raturé. $\mathfrak{Z}$ rend sa lettre en forme de 3 surmontée d'une flèche, $\mathcal{C}_E$ son C bouclé (l'ensemble des ordres circulaires sur $E$). « (inverses l'un de l'autre) » est écrit au-dessus de « circulaires » ; « arêtes » est écrit au-dessus d'un mot biffé ; « $\vec{D}_i$ » est ajouté au-dessus de « orientées » ; le $\vec{I}_s$ récrit un premier signe}

2°) Pour \struck{toute arête orientée} $\vec{a} \in \vec{A}$ de $K$, d'origine $s$ et d'extrémité $s'$, la bijection canonique
\[
(20) \qquad \psi_{\vec{a}} : \vec{\mathfrak{Z}}_s \xrightarrow{\ \sim\ } \vec{\mathfrak{Z}}_{s'}
\]
(qui donne lieu à la formule
\[
(21) \qquad \psi_{\overleftarrow{a}} = (\psi_{\vec{a}})^{-1} \text{)}
\]
\note{l'indice du premier $\psi$ de (21) est une flèche \uncertain{inversée} sur $a$ ; nous le rendons $\overleftarrow{a}$. La parenthèse ouverte avant « qui donne lieu » n'est fermée qu'au bout de (21), par nous}

Quand je dis « canonique », j'entends : provenant de la donnée de la carte $\mathcal{X}$. Inversement, si on a pu, à partir de $(K, (D_i)_{i \in I}, \omega)$, construire intrinsèquement des données (19) et (20) satisfaisant (21), \uncertain{on trouve} une immersion cellulaire combinatoire de $K$ dans une $\mathcal{X}$ (peut-être pas le plan projectif !).

\page{42}
Pour savoir que $(D_i)_{i \in I}$ y forme bien un système de ps.\ droites, il faudra d'abord savoir qu'ils se \uncertain{déduisent} de la structure de carte de la façon habituelle, ce qui s'exprime par la condition ($\forall\, s \in S$)
\[
(22) \qquad \text{si } \rho \in \vec{\mathfrak{Z}}_s \text{, et si } \nu_s \text{ est le nb de } D_i \text{ passant par } s \text{, on a } \rho^{\nu_s}(\vec{D}_i) = (-\vec{D}_i)
\]
pour toute droite orientée $\vec{D}_i$ passant par $s$.
\note{la condition est écrite sur trois lignes, sous une accolade ; « $\vec{D}_i$ » est ajouté au-dessus de « orientée »}
\marginal{\emph{NB} Si $\nu_s = 2$, alors la condition (22) implique \uncertain{déjà} la connaissance \ill{} (22) de $\mathfrak{Z}_s \subset \mathcal{C}_{\vec{I}_s}$.}

\struck{Enfin il \ill{} \ill{}} Il faudra de plus que $\mathcal{X}$ soit le plan projectif, ce qui s'exprime par la condition d'EP
\[
(23) \qquad \underbrace{c_0 - c_1}_{1 - c(K)} + c_2 = 1 \qquad \text{soit} \qquad c_2 = c(K)
\]
\note{sous $c(K)$ : « degré de connexion de $K$ », écrit au-dessous d'un mot biffé. EP : Euler-Poincaré, son abréviation}
\marginal{\ill{} la \uncertain{connaissance} de $\mathcal{X}$, i.e.\ \uncertain{l'étude} de \ill{}}
\marginal{\emph{NB} On aura a priori $c_2 \leqslant c(K)$, et la formule (23) exprime qu'on a \ill{}\ldots\ bien sûr \ill{} \ill{} i.e.\ l'\uncertain{identité}\ldots}
\note{les deux notes marginales sont écrites en oblique dans la marge gauche ; la première s'enchaîne à la ligne biffée « Enfin \ldots », la seconde est en face de (23)}

$c_0 = \operatorname{card} S$, $c_1 = \operatorname{card} A$ (nb de sommets et nb d'arêtes) ne dépendent que de $K$, tandis que $c_2$ (nb de faces) dépend bien sûr de $\mathcal{X}$.

\emph{Construction des $\vec{\mathfrak{Z}}_s$} satisfaisant (22). On va d'abord définir, pour $\nu_s \geqslant 3$
\[
(24) \qquad \mathfrak{Z}^{1}_s \subset \mathcal{C}_{I_s}
\]
\struck{\ill{}} les deux ordres circulaires opposés l'un à l'autre, sur l'ens $I_s$ des ps.\ droites passant par $s$. Soit $\Delta$ \struck{$D$} une ps.\ droite qui ne passe pas par $s$. \struck{Alors l'un des ps.\ droites \ill{} passant par $s$ est un \ill{} \ill{}}
\note{l'exposant de $\mathfrak{Z}^{1}$, ici et jusqu'à la page 45, est un petit trait crochu, \uncertain{lu} « $1$ » (peut-être un accent ou un autre signe) ; au-dessus de « ps.\ droite qui ne », une petite insertion est raturée. Les deux dernières lignes sont biffées d'un trait chacune ; en marge, en face, un renvoi oblique « \ill{} \ill{} polygonales \ill{} »}

\page{43}
\note{sa page \emph{7}, en haut à gauche}
\[
(25) \qquad I_s \xrightarrow{\ \varphi = \varphi_{s,\Delta}\ } \text{ens.\ } S_\Delta \text{ des sommets de } K \text{ sur } \Delta
\]
caractérisée par
\[
D_i \cap \Delta = \lbrace \varphi(D_i) \rbrace
\]
\note{« canonique » est écrit au-dessus de $I_s$ ; un premier symbole est raturé après $I_s$ et devant « ens. », et un numéro raturé devant la seconde formule}
Cette application est \emph{injective}, i.e.
\[
\varphi_{s,\Delta} : I_s \hookrightarrow S_\Delta
\]
Comme sur $S_\Delta$ (ens des sommets d'un polygone) il y a une structure polygonale, il y a sur $I_s$ une structure polygonale induite, soit $\mathfrak{Z}^{1}_{s,\Delta}$. On aura alors
\[
(A1) \qquad \mathfrak{Z}^{1}_{s,\Delta} \text{ est indépendant du choix de } \Delta \in I \setminus I_s \text{, soit } \mathfrak{Z}^{1}_s .
\]
\marginal{Condition portant sur les sommets $s$ tels que $\nu_s \geqslant 3$.}
\marginal{à prouver que c'est vrai quand on part d'une $\mathcal{X}$ (et qu'on trouve bien le $\mathfrak{Z}_s$ déduit de $\vec{\mathfrak{Z}}_s$\ldots) \uncertain{item} pour la \uncertain{variante} \ill{} \ill{}}
\note{la première note marginale est à gauche de (A1), la seconde, séparée du texte par un trait, court en oblique le long du bord droit jusqu'à (27)}

\struck{On} Considérons l'application
\[
(26) \qquad \vec{I}_s \longrightarrow I_s \qquad \text{de degré } 2
\]
\struck{Il faut} \add{On aura} alors
\[
(27) \qquad
\begin{cases}
\text{si } \vec{u} \in \vec{\mathfrak{Z}}_s \ (\subset \mathcal{C}_{\vec{I}_s}) \text{, } u \text{ par passage au quotient dans (26),} \\
\text{on a } u \in \mathcal{C}_{I_s} \text{, et on a } u \in \mathfrak{Z}^{1}_s .
\end{cases}
\]
\note{« on aura » est écrit au-dessus de deux mots biffés ; « $\in \mathcal{C}_{\vec{I}_s}$ » est ajouté au-dessus de $\vec{\mathfrak{Z}}_s$, et un $\varphi$ est biffé devant $\vec{u}$. Au second $u$ de la deuxième ligne, de petits signes en exposant et en indice sont illisibles}
C'est là une condition restrictive importante sur la structure polygonale $\vec{\mathfrak{Z}}_s$ qu'on veut déterminer, qui limite le choix à $2^{\nu_s - 1}$ possibilités,

\page{44}
comme on voit ainsi.
\struck{On fixe $u_0 \in \mathfrak{Z}^{1}_s$, et une origine $i_0 \in I_s$, d'où $i_1, \ldots, i_{\nu_s - 1}$ définis par $i_\nu = u_0^\nu i_0$. Soit $\omega_\nu$ l'une des deux orientations de $D_{i_\nu}$ i.e.\ l'un \ill{} \ill{} fibre \ill{} $i_\nu$ dans (26), donc $u$ \ill{} et comme}
\note{ce premier essai est encadré et barré de quatre longs traits obliques}

On \emph{fixe} $u_0 \in \mathfrak{Z}^{1}_s$, i.e.\ une orientation du polygone dont l'ens des sommets est $I_s$, pour dire que pour $i \in I_s$, on doit se donner une bijection $d_i$ de $\vec{I}_i$, la fibre de $\vec{I}_s$ en $i$, i.e.\ de l'ens des deux orientations de $D_i$, sur $\omega_j$ où $j = u_0(i)$ est le sommet suivant de $i$ --- ce qui, pour $i$ fixé, fait $2$ choix, donc en tout, a priori, $2^{\nu_s}$ choix. Mais on veut que la permutation $u$ de $\vec{I}_s$ obtenue soit une permutation circulaire, i.e.\ \struck{trans} transitive ; si on fixe $i_0 \in I_s$, cela signifie que les choix de $\alpha_{i_0}, \alpha_{i_1}, \ldots, \alpha_{i_{\nu_s - 2}}$
\[
(28) \qquad \omega_{i_0} \xrightarrow[\sim]{\ \alpha_{i_0}\ } \omega_{i_1} \xrightarrow[\sim]{\ \alpha_{i_1}\ } \omega_{i_2} \ \cdots\ \xrightarrow[\sim]{\ \alpha_{i_{\nu_s - 2}}\ } \omega_{i_{\nu_s - 1}}
\qquad (\nu = \nu_s)
\]
déterminent le choix du dernier $\alpha_{i_{\nu - 1}} : \omega_{i_{\nu - 1}} \xrightarrow{\ \sim\ } \omega_{i_\nu} = \omega_{i_0}$ (i.e.\ qu'on ait l'\emph{opposé}) par la condition qu'il soit \emph{distinct} \struck{des composés} de l'inverse du composé des bij.\ précédentes.
\struck{Dans le cas \ill{}}
\note{à gauche du texte, un dessin : un heptagone aux côtés fléchés, sommets en pointillé, avec « $u_0(i) = i+1$ » et une flèche courbe au centre. « $d_i$ » est ajouté au-dessus de « bijection », « de $\vec{I}_i$ » est \uncertain{lu}. Au-dessus de « circulaire, i.e. transitive », une insertion : « (\ill{} par $i_\alpha = u_0^\alpha(i_0)$) ». « $\nu = \nu_s$ » est écrit dans la marge gauche, en face de (28)}

\page{45}
\note{sa page \emph{8}, en haut à gauche}
Soit $\vec{\Delta}$ une ps-droite \emph{orientée} ne passant pas par $s$, donc l'orientation de $\vec{\Delta}$ en définit une de $I_s \hookrightarrow S_\Delta$, i.e.\ un générateur $u_0$, noté $u_{s,\vec{\Delta}}$, de $\mathfrak{Z}^{1}_s = \mathfrak{Z}^{1}_{s,\Delta}$. Si $\nu_s = 2$, on prendra $u_{s,\vec{\Delta}}$ la transposition dans $I_s$ (de card $2$), c'est \uncertain{ici} \ill{} \ill{} (i.e.\ \struck{\ill{}} ici $\operatorname{card} \mathfrak{Z}^{1}_s = 1$, et non $2$).
\note{« ps- » est ajouté au-dessus de « droite » ; dans « $(u_{s,\vec{\Delta}})$ » l'indice récrit un premier indice raturé}

On \struck{doit} doit définir un $\vec{u}_{s,\vec{\Delta}}$ au-dessus de $u_{s,\vec{\Delta}}$, ce qui fait (en exigeant que ce soit une permutation circulaire) ($2^{\nu_s - 1}$ \struck{\ill{}} a priori possibilités) (le compte est bon, même pour $\nu_s = 2$). On est ramené ici, pour une ps.\ droite $D_i$ passant par $s$, désignant par $D_j$ son « successeur » pour $u_{s,\vec{\Delta}}$, à décider la bijection $\omega_i \xrightarrow{\ \sim\ } \omega_j$ induite par $\vec{\Delta}$, donc pour le choix d'une orientation sur $D = D_i$, disons $\vec{D}_i$, il faut décider l'orientation correspondante de $D_j$ --- en utilisant l'orientation de $\vec{\Delta}$.
\note{au-dessus de « possibilités », une insertion : « un choix à déterminer parmi », qui renvoie au nombre $2^{\nu_s - 1}$}

\[
(29)
\]
\note{figure : un triangle formé de $\vec{\Delta}$ (horizontale, fléchée vers la droite) et de deux ps.\ droites se coupant en $s$ au sommet ; à gauche $\vec{D}_i = \vec{D}$, fléchée vers le bas, à droite $D_j = D'$ avec une flèche vers le bas à côté. Les côtés sont marqués $a$ (de $s$ à $t$), $b$ (de $s$ à $u$) et $c$ (de $t$ à $u$ sur $\Delta$) ; les pieds sont $t_i = t$ et $t_j = u$}

La figure montre que l'orientation à choisir sur $D' = D_j$ est décrite par la condition que l'arête $a$ joignant $s$ à $t = \varphi(i)$ (intersection de $D$ et $\vec{\Delta}$), dans le sens de l'orientation \struck{donnée} de $\vec{D}$, et celle $b$ joignant $s$ à $u = \varphi(j)$ (int.\ de $D'$ et $\Delta$) dans le sens de l'orientation cherchée de $D'$, et enfin l'arête $c$ joignant $t$ à $u$ dans le sens

\page{46}
de l'orientation de $\vec{\Delta}$, soit un syst.\ \struck{\ill{}} d'orientations \emph{opposé à l'orientation} \struck{des triangles} $W_{\lbrace i, j, \Delta \rbrace}$ de l'arrangement des trois ps.\ droites $\Delta$, $D$, $D'$ (ce qui caractérise l'orientation de $\vec{D}'$, parmi les deux de $D'$, en termes de celles de $\Delta$, $D$).
\marginal{i.e.\ (orientations qu'ils \uncertain{bordent}) $=$ \uncertain{triangles} \ldots}
\note{« de l' » est écrit au-dessus de « triangles » biffé ; l'indice de $W$ est \uncertain{lu} $\lbrace i, j, \Delta \rbrace$}

(A2) Cette prescription --- qui définit, pour $s$, $\vec{\Delta}$ fixés, une permutation $\vec{u}_{s,\vec{\Delta}}$ de $\vec{I}_s$ au-dessus de la permutation $u_{s,\vec{\Delta}}$ de $I_s$ --- est bien une permutation \emph{circulaire} \struck{\ill{} i.e.\ \ill{} \ill{}} i.e.\ $(\vec{u}_{s,\vec{\Delta}})^{\nu_s}$ (qui, a priori, est soit l'identité, soit la \struck{\ill{}} transposition sur chaque fibre de (26)) \struck{n'est} est la transposition.

\struck{(A3) Si $\vec{\Delta}'$ \ill{}}

\emph{NB} On aura \struck{\ill{}}
\[
(30) \qquad \vec{u}_{s,-\vec{\Delta}} = \vec{u}_{s,\vec{\Delta}}^{\,-1} , \qquad
\vec{\mathfrak{Z}}_{s,\Delta} = \lbrace \vec{u}_{s,\vec{\Delta}}, \vec{u}_{s,-\vec{\Delta}} \rbrace \subset \mathcal{C}_{\vec{I}_s}
\]
\note{après la première égalité, deux mots à droite, « pour \ill{} », sont d'une lecture incertaine}

(A3) Soit $\Delta'$ une autre ps.\ droite ne passant pas par $s$, alors
\[
\vec{\mathfrak{Z}}_{s,\Delta'} = \vec{\mathfrak{Z}}_{s,\Delta} \qquad (\text{soit } \vec{\mathfrak{Z}}_s)
\]

Bien entendu, les conditions (A1) à (A3) sont satisfaites pour un arrangement $1$-dim.\ orienté de ps.\ droites provenant d'un arrangement $2$-dimensionnel, ce qui permet donc de

\page{47}
\note{sa page \emph{9}, en haut à gauche}
construire les $\vec{\mathfrak{Z}}_s$ ($s \in S$), et \struck{les} l'application surjective
\[
(31) \qquad \vec{I} \setminus \vec{I}_s \longrightarrow \vec{\mathfrak{Z}}_s
\]
de l'ens des ps.\ dr.\ orientées ne passant pas par $s$, dans $\vec{\mathfrak{Z}}_s$, application compatible avec l'action du groupe des signes $\pm 1$.
\note{« ($s \in S$) » est écrit au-dessus de la ligne, un indice raturé après $\vec{\mathfrak{Z}}_s$ ; « surjective » est ajouté au-dessus de « application »}

Soit maintenant $\vec{a}$ une arête de $K$, joignant $s$ à $s'$, on veut définir
\[
\psi_{\vec{a}} : \vec{\mathfrak{Z}}_s \longrightarrow \vec{\mathfrak{Z}}_{s'}
\]
\struck{\ill{} ici \ill{} \ill{} besoin de} l'existence d'au moins une ps.\ droite qui ne passe ni par $s$ ni par $s'$. Considérons
\note{la ligne biffée porte au-dessus deux mots non biffés, « a priori, d'abord » (\uncertain{lecture}), qui la remplacent sans doute}
\[
\vec{J} = \vec{I} \setminus (\vec{I}_s \cup \vec{I}_{s'})
\]
\[
(32)
\]

\begin{tikzcd}[column sep=small]
& \vec{J} \arrow[dl] \arrow[dr] & \\
\vec{\mathfrak{Z}}_s \arrow[rr, dashed, "\psi_{\vec{a}}"] & & \vec{\mathfrak{Z}}_{s'}
\end{tikzcd}

on obtient deux quotients de $\vec{J}$ --- on a donc la condition que c'est le même, i.e.

(A4) si $\vec{\Delta}$, $\vec{\Delta}'$ sont deux ps.\ droites orientées ne passant pas par les extrémités $s$, $s'$ de l'arête $a$, qui induisent la même orientation en $s$, elles induisent aussi la même orientation en $s'$.

\page{48}
Reste à traiter le cas (très particulier !) où $\nexists$ ps.\ droite \struck{$\Delta$} qui ne passe par $s$ ni par $s'$.
Soit $\Theta$ la ps.\ droite déterminée par $a$, considérons donc une ps.\ droite $D$ ($\neq \Theta$) passant par $s$, une $D'$ ($\neq \Theta$) passant par $s'$, \struck{soit} où donc $D \neq D'$, soit $t$ l'intersection de $D$ et $D'$. Soit, \struck{$b$ un des deux} $K' = K_{\lbrace \Theta, D, D' \rbrace}$,
\note{le « $\nexists$ » est son $\exists$ barré ; « ($\neq \Theta$) » est ajouté deux fois au-dessus de la ligne}

\[
(33)
\]
\note{figure : un triangle de sommets $s$, $s'$ (sur $\Theta$, horizontale, avec $\vec{a}$ fléchée de $s$ vers $s'$) et $t$ au sommet ; le côté $b$ de $s$ à $t$ et le côté $b'$ de $t$ à $s'$ sont fléchés dans le sens $s \to t \to s'$, plusieurs petites flèches le long de chacun. Une double flèche verticale, dans la marge, descend vers le texte qui suit}

soit $b$ une des arêtes qui joignent $s$ à $t$, et prenons l'unique arête $b'$ de $K'$ qui joigne $t$ à $s'$, et telle que $\lbrace a, b, b' \rbrace$ soit compatible avec l'orientation de $K'$ ! Alors $b$ et $b'$ ne sont pas des arêtes de $K$, mais ils \uncertain{sont} \ill{} des arêtes de $K$, \uncertain{se} décomposent en succession d'arêtes de $K$.
\marginal{\ill{} \ill{} : \uncertain{comme} $\vec{\psi}_{\vec{a}} : \vec{\mathfrak{Z}}_s \to \vec{\mathfrak{Z}}_{s'}$, comme on connaît, quand on a un ordre circulaire $\in \mathfrak{Z}_s$, la restriction de l'ordre \ill{} sur \struck{\ill{}} la partie de $\vec{I}_s$, \ill{} deux droites $\Theta$, $D'$\ldots}
\note{la note marginale est écrite en oblique, en lignes serrées, dans la marge gauche en face de ces lignes}

On suppose que $K$ n'est pas réduit à \struck{une sécante} (famille de) $n \geqslant 3$ ps.\ droites concourantes et une seule sécante (ce qui signifie que les $n$ ps.\ droites passant par le sommet privilégié $s$ sont toutes des bigônes). Dans ce cas, la ps.\ dr.\ $\Theta$ est l'\emph{unique} ps.\ droite qui ne soit pas un bigône, donc $D$ et $D'$ n'en sont pas \ill{} de $K$, \ill{} définir les $\psi_{\vec{c}}$ pour des arêtes \uncertain{portées} par $D$, ou par $D'$. On définit donc
\marginal{\emph{NB} Obs.\ que $K$ \struck{\ill{}} \ill{} ($4$ ps.\ \ill{}) \ill{} $3$ \ill{} ps.\ droites \ill{} (cas \ill{} déjà traité !)}
\note{« réduit à » est écrit au-dessus de « une sécante » biffé ; « ps.\ dr. » récrit « droite » biffé. La note marginale d'en bas, oblique, commence sous l'autre et déborde sur le texte ; on n'en lit que des fragments. Les deux dernières lignes de la page sont d'une lecture très incertaine}

\page{49}
\note{sa page \emph{10}, en haut à gauche}
\[
(34) \qquad \vec{\psi}_{\vec{a}} = \psi_{\vec{b}'} \, \psi_{\vec{b}}
\]
où $\vec{b}$ est $b$ orienté de $s$ vers $t$, $\vec{b}'$ est $b'$ orienté de $t$ vers $s'$, et $\psi_{\vec{b}'}$ et $\psi_{\vec{b}}$ sont définis comme composés des $\psi_{\vec{c}}$ correspondants aux arêtes de $K$ composantes des chemins $\overrightarrow{(s,t)}$ et $\overrightarrow{(t,s')}$.
\note{ici les $\psi$ sont écrits petits, presque comme des « v » ; « de $K$ » est ajouté au-dessus de « composantes »}

Il faudrait, dans cette optique, pour une (A5) --- savoir que ce $\psi_{\vec{a}}$ ne dépend pas des choix arbitraires qui ont été faits. Mais le cas où $K$ comporte une ps.\ droite qui est un bigône \struck{\ill{} \ill{}} devrait plutôt être traité à part, avec des moyens plus directement adaptés à cette situation. Je n'ai donc pas vérifié que la proposition énoncée soit correcte, sans exclure le cas où $K$ est formé de $n$ pseudo-droites concourantes, coupées par une seule autre ps.\ droite.
\marginal{mais \uncertain{si} l'\ill{} \ill{} de la \ill{} (v.\ p.\ précédente)}
\note{« devrait plutôt être » est écrit au-dessus d'une fin de ligne biffée. La note marginale, oblique, est rattachée par un trait vertical aux lignes « Je n'ai donc pas vérifié \ldots\ ps.\ droite »}

J'ai maintenant envie de montrer que dans le cas d'un $K$\ldots\ provenant d'une $\mathcal{X}$, son orientation peut se décrire canoniquement en termes de $(K, (D_i))$ seulement.

\page{50}
Soit donc un « triangle » subordonné à $(K, (D_i)_{i \in I})$, \struck{\ill{}} défini par $J \in \mathfrak{P}_3(I)$, formé de trois ps.\ droites \uncertain{concourantes}. On veut décrire \struck{\ill{}} l'orientation « canonique » de $K_J$.
\note{on attendrait « non concourantes » pour un triangle ; nous ne lisons pas de « non » sur la page}

(1er cas) Il existe une \struck{\ill{}} ps.\ droite $\Delta$, qui ne passe par aucun des trois sommets de $K_J$ --- i.e.\ telle que les systèmes $K_{J \cup \lbrace i \rbrace}$ des $4$ ps.\ droites soit formé d'un syst.\ de quatre ps.\ droites $3$ à $3$ non concourantes. On utilise
\note{un signe raturé devant « (1er cas) », et devant $\Delta$}

\emph{Proposition} Le \struck{catégorie des arrangements} foncteur \struck{\ill{} \ill{}} \uncertain{évident}
\[
\begin{pmatrix} \text{arrangements simples} \\ \text{de } 4 \text{ ps.\ droites} \end{pmatrix}
\longrightarrow
\begin{pmatrix} \text{arr.\ } 1\text{-dimensionnels} \\ \text{simples de } 4 \\ \text{ps-droites} \end{pmatrix}
\]
est une équivalence de catégories.
\note{« foncteur » et les mots qui suivent sont écrits au-dessus de la ligne biffée. Au-dessus de la formule, deux lignes commencées sont biffées : « 1 -- \ill{} », « \ill{} de $4$ ps.\ droites », « \ill{} \ill{} »}

(On le démontre en comparant les deux \uncertain{groupoïdes}, à celui des ens.\ \struck{finis} à $4$ éléments)

Donc, comme $K_{J \cup \lbrace i \rbrace}$ provient canoniquement d'un arrangement $2$-dimensionnel, $K_J$ hérite d'une orientation, soit $\omega_{J,i}$, et on aura

\page{51}
\note{sa page \emph{11}, en haut à gauche}
(A5) l'orientation $\omega_{J,i}$ de $K_J$ est indépendante du choix de $i$.

\emph{« 2ème cas »} Toute ps.\ droite de $K$ passe par un des trois sommets de $K_J$.
Comme on suppose qu'il n'y a pas de bigône dans $K$, on trouve d'ailleurs que par chaque sommet, il passe au moins une $D_i$ avec $i \in I \setminus J$.
\marginal{Si on désigne par $\alpha$, $\beta$, $\gamma$ les intersections \ill{} \ill{} des trois ps.\ droites $D_i$ \ill{} de $K_J$, et par $\partial\alpha$, $\partial\beta$, $\partial\gamma$ \ill{} \ill{} arêtes de $K_I$ qui \ill{} \ill{} l'orientation \ill{} $\partial\alpha$, $\partial\beta$, $\partial\gamma$ \ill{}}
\note{la note marginale est écrite en oblique dans le coin supérieur gauche, autour de « (A5) » et de « 2ème cas ». Au-dessous, un dessin : un triangle de sommets $a$, $b$, $c$, les trois côtés prolongés}

Supposons d'abord qu'on peut trouver des droites $D_a$, $D_b$, $D_c$ passant par les sommets $a$, $b$, $c$, qui ne soient pas concourantes. On trouve une configuration, associée canoniquement à l'ens $J$ à trois éléments, où comme sommets les $a_i$ ($i \in J$) de $K_J$, les intersections $\alpha_i$ ($i \in J$) des $D_{a_i}$ avec $D_i$ (\emph{NB} $a_i$ est le sommet de $K_J$ opposé à $D_i$) et les $a'_i$ ($i \in J$) intersection de $D_{a_j}$ et $D_{a_k}$ (où $J = \lbrace i, j, k \rbrace$). \struck{\ill{}} La figure montre comment les sommets se répartissent sur les $6$ ps.\ droites $D_i$, $D_{a_i}$, et \struck{la configuration est} (avec les $6$ ps.\ droites) est complètement décrite, quand on a décrit les structures polygonales
\note{figure, à gauche : six droites, les côtés du triangle $a$, $b$, $c$ et trois droites $D_a$, $D_b$, $D_c$ issues des sommets ; points marqués $a$, $b$, $c$, $a'$, $b'$, $c'$ à l'intérieur, $\alpha$, $\beta$, $\gamma$ sur les côtés. « $9$ » est écrit en marge devant « sommets ». Au-dessus de « répartissent », un indice \uncertain{$D_{J'}$} ; au-dessus de « la configuration est » biffé, « correspondant », \uncertain{lu}}

\page{52}
sur chacun des quatre sommets $a_i$, $\alpha_i$, $a'_j$, $a'_k$ \struck{\ill{}} se trouvant sur une $D_{a_i}$, i.e.\ en \struck{\ill{}} précisant lequel, parmi les trois sommets $\alpha_i$, $a'_j$, $a'_k$ distincts de $a_i$, est \emph{opposé} à $a_i$. On voit que si la figure $1$-dimensionnelle provient d'un arrangement $2$-dimensionnel (cas I), il y a un ordre circulaire sur $J$ (unique), tel que $a_i$ soit opposé, sur $D_{a_i}$, à $a'_{i+1}$. (\emph{NB} Dans le cas de la figure choisie, cet ordre circulaire est $abc$).
\marginal{Il y a une \uncertain{autre} possibilité : pour $i \in J$, $a_i$ est opposé à $\alpha_i$ (cas II)}
\note{« chacun » est écrit au-dessus d'un mot biffé ; « (cas I) » est ajouté au-dessous de « $2$-dimensionnel ». La note marginale, oblique, porte sa parenthèse « (cas II) » au bas de la marge}

\struck{On trouve les configurations} \struck{On voit que \ill{}} \ill{} \ill{} \struck{$2$-dimensionnelles} $K_{J'}$, \struck{\ill{}} et $K_J$, provenant d'un arr.\ $2$-dimensionnel, donc, \struck{les} désignant par $\lambda_i$ ($i \in J$) \struck{les} l'arête de $K_J$ \struck{qui \ill{}} portée par $D_i$ qui contient $\alpha_i$, les $\lambda_i$ ($i \in J$) forment un syst.\ d'arêtes compatible avec l'orientation dans le cas I, et non compatible dans le cas II.

(A6) \struck{L'orientation} Dans le cas d'un système $J \in \mathfrak{P}_3(I)$ non concourant, de trois ps.\ droites, tel que toute autre ps.\ droite passe par un des trois sommets $a_i$ de $K_J$, et lorsqu'on peut trouver pour chaque $a_i$ une $D_{a_i}$ qui passe par $a_i$ et

\page{53}
\note{sa page \emph{12}, en haut à gauche}
\struck{distincts de $\ell$} $\in I \setminus J$, de façon que les trois $D_{a_i}$ soient non concourantes, on exige

(1°) que l'arr.\ $1$-dim.\ \struck{configuration} $K_{J'}$ des $6$ droites obtenu soit du cas (I) ou (II) ci-dessus, et

(2°) que l'orientation de $K_J$ définie par la prescription précédente ne dépende pas des choix de $J'$.

Il reste enfin à regarder le cas où \ill{} les systèmes de $D_{a_i}$ ($D_{a_i} \in I \setminus J$ passant par $a_i$), celles-ci sont concourantes. Mais cela implique que \struck{l'arrangement est} l'arrangement \ill{} $\operatorname{card} I = 6$, et que \ill{} standard formé de $7 = 4 + 3$ sommets (\uncertain{qui sont des triangles}), et $6$ ps.\ droites, les six ps.\ droites étant les ps.\ droites qui relient deux à deux quatre sommets triples, les trois points d'intersection supplémentaires étant les intersections des ps.\ droites « opposées ». La situation formée de ps.\ droites $1$-dimensionnelle (un complexe combinatoire \ill{} par \ill{} triangles) est entièrement décrite par l'ens des quatre sommets triples. La situation $2$-dimensionnelle correspondante ($=$ \uncertain{groupe d'automorphismes} $\mathfrak{S}_4$) est aussi standard, \ill{} la description suivante : On trouve
\marginal{\ill{} \uncertain{pour} \ill{} $D_i$ ($i \in I \setminus J$) \ill{} que chacune \ill{} par un des trois sommets}
\note{la note marginale, oblique, est rattachée au mot « où » ; « isomorphe à : » est écrit au-dessus de la fin de la ligne biffée. À gauche, un dessin : un triangle $a$, $b$, $c$ avec ses trois médianes concourantes en $d$ et une sixième droite, les pieds marqués $\alpha$, $\beta$, $\gamma$ ; les sommets $a$, $b$, $c$ sont des points noirs. « $7$ » est son 7 barré. L'arrangement de six ps.\ droites à quatre sommets triples est, semble-t-il, celui que le lot 2 (p. 24) note $\mathrm{III}_6$}

\page{54}
Soient $a$, $b$, $c$ les sommets de $K_J$, $\alpha$, $\beta$, $\gamma$ les sommets « opposés » dans $K_{J'} = K_I$, intersection de $D_a$ avec la ps.\ dr.\ joignant $b$ à $c$ etc., soient $\lambda_a$, $\lambda_\beta$, $\lambda_c$ les arêtes de $K_I$ contenant $\alpha$, $\beta$, $\gamma$, alors $\lbrace \lambda_a, \lambda_b, \lambda_c \rbrace$ est compatible avec l'orientation.
\note{les indices des $\lambda$ sont écrits tels qu'ils se lisent : $\lambda_\beta$ dans la première énumération, $\lambda_b$ dans l'accolade}

En résumé

\emph{Théorème} Le foncteur canonique du groupoïde des arrangements de \struck{pseudodroites} sans bigônes, vers celui des arrangements $1$-dimensionnels de pseudo-droites (sans bigônes) est pleinement fidèle.
\note{« celui » est écrit au-dessus de « des », rattaché par un trait}

Les conditions (A1) à (A6) donnent en principe une voie pour déterminer l'image essentielle --- mais il faudrait d'abord simplifier les conditions en question. En principe, il faudrait encore rajouter la relation d'EP (23) --- mais je soupçonne que cette dernière est automatique\ldots

\page{56}
\note{petite carte à l'italienne, dessin à l'encre sans texte suivi : un arrangement de droites à cinq sommets multiples (points noirs), deux verticales, deux horizontales et des obliques en éventail. Chaque région porte le nombre de ses côtés, $3$, $4$ ou $5$ (une seule région $5$, en bas au centre) ; les régions à l'infini, en haut, sont marquées « id » (ou « i$\delta$ »). À droite, le décompte}
($+\ 9$ triangles $+\ 1$ carré)

$25$ triangles

$5$ carrés

$1$ pentagone

$31$ régions
\note{le « $25$ » récrit un premier chiffre ; après « $31$ » un mot est biffé. Le décompte « $31$ régions » est isolé par un double trait vertical}

\page{57}
\note{petite carte à l'italienne, dessin à l'encre sans texte : six droites, deux verticales, deux presque horizontales et deux obliques, qui forment au centre un petit triangle entouré de triangles ; les régions sont marquées $3$ ou $4$, celles du haut « id » (ou « i$\delta$ »). À droite, quatre points isolés}

\page{58}
\note{petite carte à l'italienne : un arrangement d'une dizaine de droites, sans sommet marqué, dont certaines extrémités portent les étiquettes $5_1$, $5_2$, $5_4$, $5_5$, $10_3$, $1_6$, et une région centrale « $4$ ». À droite, un décompte et une légende}
\[
\begin{array}{r}
5 \\
15 = 5 + 10 \\
10 = 5 + 5 \\
1 \\
\hline
31
\end{array}
\]
$1_6 \equiv{}$ icosaèdre

$5_1 \equiv{}$ hexagone

$5_4$, $5_2 \equiv{}$ triangle bordé

$5_5$, $10_3 \equiv{}$ carré \uncertain{gammé}
\note{le « $10$ » de la troisième ligne récrit un premier chiffre}

\section{Système de Pseudodroites}
\note{titre souligné, de sa main, en tête de la page 60, écrite au crayon et non numérotée}

\page{60}
\emph{Définition} : c'est une carte (topologique ou combinatoire) dont la surface sous-jacente $X$ est un plan projectif${}^{*}$ réel, dont les ordres des sommets sont pairs $\geqslant 4$ (d'où la division de l'ensemble des arêtes en classes d'équivalence, appelées les « pseudodroites » de la carte) telles que deux pseudodroites se rencontrent en un point et un seul.
\marginal{${}^{*}$ i.e.\ (compacte) non orientable \struck{et} de car.\ EP égale à $1$\ldots}
\note{la note, appelée par l'astérisque, est écrite verticalement dans la marge gauche}

\emph{Premières conséquences} : le nb $n$ des pseudodroites est $\geqslant 2$. L'immersion de $\mathbf{S}^1$ dans $\mathbf{P}^2_{\mathbf{R}}$ correspondant à une ps.\ dr.\ est (des deux classes d'isotopie possibles) celle qui est \uncertain{non} homotope à $0$, i.e.\ le long d'une ps.\ dr.\ $X$ est non orientable, i.e.\ la situation est la même que pour une \emph{droite} (projective) dans le plan projectif. On montre que pour des ps dr « en position générale » (sommets tous d'ordre $4$) et en nb $n \leqslant 6$, la situation peut être réalisée par un système de vraies droites dans le plan projectif réel.
\note{le « $n$ » est ajouté au-dessus de « nb », après un signe raturé ; le « $0$ » de « homotope à $0$ » récrit un premier signe}

Soit $D$ une pseudodroite dans $X$, $\widetilde{X}$ le revêtement des orientations de $X$ ($\widetilde{X} \simeq \mathbf{S}^2 \simeq \mathbf{P}^1_{\mathbf{C}}$), $\widetilde{D}$ l'image inverse de $D$ dans $X$. Alors $\widetilde{X} \setminus \widetilde{D}$ a deux comp.\ connexes (ou « hémisphères ») en corr.\ 1-1 can.\ avec l'ens.\ des deux orientations de $D$.
\note{« dans $X$ » est écrit tel quel, où l'on attendrait $\widetilde{X}$}

Notons que \add{sur} l'ens.\ des sommets se trouvant
\note{« sur » est ajouté au-dessus de la ligne ; la phrase continue au lot suivant}

\end{document}
